\DocumentMetadata{
  pdfversion=2.0,
  pdfstandard=ua-2,
  lang=en-US,
  tagging=on,
  tagging-setup={math/setup=mathml-SE,tabsorder=structure}
}
\documentclass[11pt,letterpaper]{article}
\usepackage[margin=0.68in,includefoot]{geometry}
\usepackage{newcomputermodern}
\usepackage{amsmath,amscd,mathtools}
\usepackage{tikz,adjustbox}
\providecommand{\square}{\mdlgwhtsquare}
\usepackage[hidelinks]{hyperref}
\hypersetup{
  pdftitle={Logic PhD Exam, May 2000},
  pdfauthor={University of Florida Department of Mathematics},
  pdfsubject={Graduate mathematics examination},
  pdfdisplaydoctitle=true,
  bookmarksopen=true
}
\setcounter{secnumdepth}{0}
\setlength{\parindent}{0pt}
\setlength{\parskip}{0.28em}
\setlength{\emergencystretch}{3em}
\setlength{\leftmargini}{2.15em}
\setlength{\leftmarginii}{2.65em}
\setlength{\leftmarginiii}{3em}
\pagestyle{empty}
\raggedbottom
\allowdisplaybreaks
\NewTaggingSocketPlug{block/list/label}{examproblem}{%
  \tagstructbegin{tag=Lbl}\tagstructend%
  \tagstructbegin{tag=\LBody}%
  \tagstructbegin{tag=\ProblemHeadingTag,title-o={\ProblemHeadingText}}%
  \tagmcbegin{tag=\ProblemHeadingTag,actualtext={\ProblemHeadingText}}%
  #2%
  \tagmcend%
  \tagstructend%
}
\newenvironment{examproblems}[2]{%
  \def\ProblemHeadingTag{#2}%
  \AssignTaggingSocketPlug{block/list/label}{examproblem}%
  \begin{enumerate}%
  \setcounter{enumi}{#1}%
  \renewcommand{\labelenumi}{\textbf{\theenumi.}}%
  \setlength{\topsep}{0.35em}%
  \setlength{\partopsep}{0pt}%
  \setlength{\itemsep}{0.48em}%
  \setlength{\parsep}{0.18em}%
}{\end{enumerate}}
\newenvironment{examsubparts}{%
  \AssignTaggingSocketPlug{block/list/label}{default}%
  \begin{enumerate}%
  \setlength{\topsep}{0.18em}%
  \setlength{\partopsep}{0pt}%
  \setlength{\itemsep}{0.16em}%
  \setlength{\parsep}{0.1em}%
}{\end{enumerate}}
\newenvironment{examinstructions}{%
  \par\begingroup\itshape
}{%
  \par\endgroup\vspace{0.18em}
}
\makeatletter
\renewcommand\subsection{\@startsection{subsection}{2}{0pt}%
  {-1.25ex plus -.25ex minus -.1ex}%
  {0.55ex plus .1ex}%
  {\normalfont\large\bfseries}}
\renewcommand{\@maketitle}{%
  \newpage\null\vskip -1em%
  {\footnotesize\bfseries
    \href{https://www.ufl.edu/}{University of Florida}, \href{https://math.ufl.edu/}{Department of Mathematics}\par}%
  \vskip -0.5em%
  \tagstructbegin{tag=H1,title={Logic PhD Exam, May 2000}}%
  \tagpdfparaOff%
  \tagmcbegin{tag=H1}%
    {\LARGE\bfseries\href{https://gma.math.ufl.edu/exams/logic-phd/}{Logic PhD Exam}, May 2000\par}%
  \tagmcend%
  \tagpdfparaOn%
  \tagstructend%
  \vskip 0.25em%
  \tagmcbegin{artifact}%
    \hrule height 1pt%
  \tagmcend%
  \vskip 1em%
}
\makeatother
\title{Logic PhD Exam, May 2000}
\author{}
\date{}
\begin{document}
\maketitle
\thispagestyle{empty}
\begin{examinstructions}
Answer 7 questions, including at least one from each of the four segments.
\end{examinstructions}
\subsection{1. General Logic}
\begin{examproblems}{0}{H3}
\def\ProblemHeadingText{Problem 1.}
\item
State Gödel’s incompleteness theorem, and explain the role of Gödel numbering in its proof.
\def\ProblemHeadingText{Problem 2.}
\item
Prove that there is no finite set~\(\Gamma\) of sentences in the language of groups such that a group~\(G\) satisfies all the sentences in~\(\Gamma\) if and only if~\(G\) is not torsion free. [The group~\(G\) is torsion free if \(g^n \ne e\) for all \(g \ne e\) in~\(G\).]
\def\ProblemHeadingText{Problem 3.}
\item
Show that any complete, finitely axiomatizable first-order theory is decidable.
\def\ProblemHeadingText{Problem 4.}
\item
Give an example of two models which are elementarily equivalent, but not isomorphic.
\end{examproblems}
\subsection{2. Model Theory}
\begin{examproblems}{0}{H3}
\def\ProblemHeadingText{Problem 1.}
\item
State and sketch a proof of the downward Löwenheim–Skolem theorem.
\def\ProblemHeadingText{Problem 2.}
\item
State the omitting types theorem, and use it to prove that any~\(\omega\)-complete theory for arithmetic has an~\(\omega\)-model.
\def\ProblemHeadingText{Problem 3.}
\item
State and sketch a proof of the Łoś theorem about ultraproducts of models.
\end{examproblems}
\subsection{3. Set Theory}
\begin{examproblems}{0}{H3}
\def\ProblemHeadingText{Problem 1.}
\item
Prove that a countable union of countable sets is countable.
\def\ProblemHeadingText{Problem 2.}
\item
Sketch the proof that \(L \models \mathrm{GCH}\).
\def\ProblemHeadingText{Problem 3.}
\item
Prove that if~\(\kappa\) is a regular, uncountable cardinal then any diagonal intersection of closed and unbounded subsets of~\(\kappa\) is closed and unbounded.
\end{examproblems}
\subsection{4. Recursion Theory}
\begin{examproblems}{0}{H3}
\def\ProblemHeadingText{Problem 1.}
\item
Show that a function \(h: \omega \to \omega\) is (partial) recursive if and only if its graph is recursively enumerable; and that if~\(h\) is total then its graph is recursive.
\def\ProblemHeadingText{Problem 2.}
\item
State the halting problem, and prove that it is undecidable.
\def\ProblemHeadingText{Problem 3.}
\item
Explain why, for any \(A \subset \omega\), there are at most countably many sets~\(B\) such that \(B <_T A\).
\end{examproblems}
\end{document}
